\section{Results}

Unless stated otherwise, reliability is ICC(A,1) (two-way, absolute agreement, single measure) between the two sessions of the first round (R1 S1--S2, about 30 min apart), the unit is the recording-level $A$ of each session, and intervals are 95\% participant-bootstrap percentile intervals. The intervals are descriptive: they are not adjusted for the number of constructions, floors, cohorts and contrasts compared, and the admitted constructions were defined after the results were seen (Methods).

\subsection{Constructing the interocular reading-event contrast}

The reading task of GazeBaseVR comprises 1,004 binocular recordings at 250\,Hz from 407 participants: 814 recordings from the two R1 sessions of all 407 participants, 70 from 35 participants in R2 and 120 from 60 participants in R3~\citep{lohr2023}. In each recording we classified saccades operationally by horizontal amplitude alone, as leftward return sweeps (RS; horizontal amplitude below $-15^\circ$), the class of saccades that take the eyes from the end of one line to the start of the next~\citep{parker2019vr,slattery2019}, or rightward progressive saccades ($0.5$--$10^\circ$) within a line. For each saccade, $l = \log(\PV_L/\PV_R)$ is the log ratio of left-eye to right-eye peak velocity, and the marker is
\begin{equation}
A = \tfrac12\left[\med_{\mathrm{RS}}(l) - \med_{\mathrm{prog}}(l)\right]
\label{eq:A}
\end{equation}
per recording (\cref{fig:schematic}). Because RS and progressive saccades move in opposite horizontal directions, the classic finding that the abducting eye moves faster~\citep{collewijn1988} predicts a positive $A$. Binocular coordination during return sweeps has itself been studied in skilled adult readers~\citep{parker2019jov}. Because the two classes also differ in amplitude and in their role in reading, $A$ is treated here as an exploratory interocular reading-event contrast. $A$ is dimensionless; in the joint cohort defined below, the participant mean of the two R1 sessions' $A$ under admitted events had a mean of .038 and an SD of .040, and was positive in 83\% of participants. Because $\exp(2A)$ is the ratio of the exponentiated class medians of $l$, a value of .038 corresponds to a typical return-sweep left/right peak-velocity ratio about 8\% higher than the progressive one.

% === Fig 1: schematic (FigEdit; source figures/fig1_schematic.svg) ===
\begin{figure}[t]
\centering
\includegraphics[width=\textwidth]{figures/fig1_schematic.pdf}
\caption{\textbf{Event construction for the interocular reading-event contrast $A$.}
Left- and right-eye gaze traces are converted to per-eye angular velocity. Saccades are detected either on the mean of the two per-eye speed
magnitudes (mean-speed detector) or separately in each eye (per-eye detector). The admission gate keeps a mean-speed event only when it overlaps a
matched left/right per-eye event pair, whether or not that pair passes its own exclusion rules; detection windows, peak samples and saccade classes of admitted events are unchanged. Centring is an optional
branch that keeps events whose log left/right peak-velocity ratio $l$ lies within $\log 2$ of the recording median. $A$ is half the difference between
the median $l$ of return sweeps and of progressive saccades in a recording. A recording enters only when it meets the event floor; at floor 5 it
must contain at least 5 return sweeps and 20 progressive saccades. Inset: R1 S1--S2 ICC(A,1) of unfiltered mean-speed events, of admitted events and of mean-speed events with a post hoc 1,000\textdegree/s peak-velocity ceiling,
in the joint cohort, the participants whose two R1 sessions were eligible under all six constructions ($n$ = 376, floor 5), and the first two again with the
single most influential participant omitted ($n$ = 375; Results).
Detector thresholds and the six-construction lookup are given in Methods.}
\label{fig:schematic}
\end{figure}
% inset values: unfiltered mean-speed .374, admitted .693 (pw_p1_joint_icc.csv), ceiling .711 (rv_v7_ceiling_icc.csv), omitted .639 vs .693 (rv_v4_loo_summary.csv)


Six constructions of $A$ were compared (Methods). Mean-speed detection thresholds the mean of the two per-eye speed magnitudes and reads each eye's peak within the common event window. Per-eye detection thresholds each eye separately and pairs the two eyes' events one to one. Admission keeps a mean-speed event only when it overlaps a matched per-eye pair. It removes events but leaves the windows, peaks and classes of the remaining events unchanged. Centring is an optional abstention rule that keeps an event only when its $l$ lies within $\log 2$ of the recording median. Crossing the three event sets (mean-speed, admitted, per-eye) with unfiltered or centred gives the six constructions: mean\_u, mean\_c, adm\_u, adm\_c, pereye\_u and pereye\_c.

The mean-speed detector returned 255,303 events, of which 138,642 entered the RS and progressive pools after duration and missingness exclusions (15,374 RS, 123,268 progressive); centring kept 126,386 of them (.912). The per-eye detector returned 196,983 matched events, 126,237 in the pools, of which centring kept .977 (Supplementary Table S3). Of the mean-speed pool events, 11,535 of 123,268 progressive saccades (9.4\%) and 417 of 15,374 RS (2.7\%) overlapped no matched per-eye pair. The interocular log ratios of these unsupported events were much further from equality ($l = 0$): the median of $|l|$ was 1.681 for unsupported progressive saccades and .141 for supported ones, and the 95th percentile of $|l|$ for unsupported RS was 5.462 (Supplementary Table S4). At floor 5 (at least 5 RS and 20 progressive saccades per recording), 379 participants had both R1 sessions eligible with unfiltered mean-speed events and 377 with centred mean-speed events; 376 were eligible under all six constructions. We call these 376 participants the joint cohort.

\subsection{Short-term reliability depended on event construction and recording eligibility}

In the joint cohort ($n$ = 376, floor 5, R1 S1--S2), admitting only mean-speed events that overlapped a matched per-eye pair raised the ICC(A,1) of $A$ from .374 [.202, .719] for unfiltered mean-speed events (mean\_u) to .693 [.612, .761] for admitted events (adm\_u), a path-specific admission contrast mean\_u $\to$ adm\_u of $\Delta$ICC = .319 [.003, .501], with windows, peaks and classes unchanged (\cref{fig:icc}a,b). This contrast therefore reflects the change in the event set alone. In the same cohort the other four constructions had ICC(A,1) of .714 [.654, .765] for pereye\_u, .726 [.669, .774] for mean\_c, .727 [.673, .774] for adm\_c and .717 [.657, .766] for pereye\_c. Admission restriction reproduced most of the observed point increase. On the joint cohort, the point increase from mean\_u was .319 to adm\_u, against .353 to adm\_c and .342 to pereye\_c (differences computed from the unrounded joint-cohort ICC estimates). The path-specific contrast mean\_u $\to$ mean\_c was .351 [.016, .524]. The comparison of unfiltered with centred mean-speed events was part of the frozen specification; its evaluation in this joint cohort, like the admission contrast, was added afterwards (Methods).

% === Fig 2: repeatability by construction (gen_fig2_icc.py) ===
\begin{figure}[t]
\centering
\includegraphics[width=\textwidth]{figures/fig2_icc.pdf}
\caption{\textbf{Short-term repeatability of $A$ under each event construction, with one participant omitted, under recording eligibility and under detector settings.}
\textbf{a}, ICC(A,1) between R1 sessions 1 and 2 with 95\% participant-bootstrap CI (5,000 draws) for the six constructions, in the joint
cohort of participants eligible under all six ($n$ = 376, floor 5); the unit is the recording-level $A$ of each session. Open squares, unfiltered; filled circles, centred;
orange, mean-speed events; blue, admitted mean-speed events; green, per-eye events. $\ddagger$, post hoc comparator not in the frozen plan: mean-speed events after removing those with
either per-eye peak velocity above 1,000\textdegree/s (same cohort; .711 [.642, .769]).
\textbf{b}, Path-specific contrasts $\Delta$ICC (end $-$ start; u, unfiltered; c, centred) with 95\% CI from paired participant-bootstrap draws. Filled diamonds, all participants ($n$ = 376);
open circles, the same contrasts with the participant who contributed most to the session differences under unfiltered mean-speed events omitted ($n$ = 375). The contrasts share endpoints and are not
an additive partition of reliability. With windows, peaks and classes unchanged, admission raised ICC from .374 to .693,
$\Delta$ICC = .319 [.003, .501], or .054 [$-$.007, .129] with that participant omitted;
the ceiling comparator reproduced the increase, and admitted minus ceiling was $-$.018 [$-$.090, .039]. \textbf{c}, Mean-speed detector under different eligibility, floors 5 and 10:
each construction's own cohort, the common cohort eligible both unfiltered and centred, and a count-matched random-exclusion comparator ($\dagger$)
that keeps, in each participant, session and class, the number of events retained by centring (mean of 200 repetitions; its dashed bars without caps show the
2.5--97.5 percentile range of repetitions, not a CI). Dotted lines separate the own-cohort row from the common-cohort comparisons. Right-hand text: contrasts with 95\% CI in the common cohort ($n$ = 377 at floor 5;
$n$ = 342 at floor 10). \textbf{d}, Post hoc detector-setting sensitivity: all R1 recordings re-extracted with one detector setting changed at a time (velocity threshold, Savitzky--Golay window, maximum
left/right pairing gap in samples), each in its own joint cohort (right-hand $n$; floor 5; point estimates). Unfiltered mean-speed ICC ranged from .226 to .703,
admitted from .661 to .706 and ceiling from .697 to .736; re-extraction with the frozen settings reproduced the analysed events exactly.}
\label{fig:icc}
\end{figure}


The steps after admission produced smaller point changes with uncertain increments. In the joint cohort ($n$ = 376, floor 5, R1 S1--S2), replacing admitted mean-speed events with matched per-eye events (adm\_u $\to$ pereye\_u) changed the ICC by +.021 [$-$.019, .083]. This comparison changes event membership and measurement windows and can also change class assignment. Centring the admitted events (adm\_u $\to$ adm\_c) changed it by +.034 [$-$.008, .100]. Centring the per-eye events (pereye\_u $\to$ pereye\_c) changed it by +.002 [$-$.014, .018], and the centred per-eye and mean-speed constructions differed by $-$.009 [$-$.025, .007] (mean\_c $\to$ pereye\_c). All contrasts use the same participants and bootstrap draws. Contrasts telescope numerically along a specified path, but they do not identify separate, order-independent contributions of admission, detection and centring.

Which recordings were eligible also moved the estimate, separately from admission. For unfiltered mean-speed events in their own cohort ($n$ = 379, floor 5, R1 S1--S2), ICC(A,1) was .170 [$-$.004, .693]. Restricting the same construction to the participants who were also eligible with centred events ($n$ = 377) gave .375 [.204, .721] (\cref{fig:icc}c). The two participants removed were eligible only unfiltered (Supplementary Table S5). One of them was the largest contributor to the own-cohort summed squared session differences (53\%; see the influence analysis below), so this change mainly reflects the exclusion of one participant with one extreme session. This is a change of cohort, not of construction, and it is not part of the joint-cohort contrasts above.

To ask whether centring helped because of which events it removed or only because of how many, we compared it with a count-matched random-exclusion comparator. In each participant, session and class, it kept a random subset of unfiltered mean-speed events as large as the subset kept by centring (200 repetitions; Methods). At floor 5, in the participants eligible both unfiltered and centred ($n$ = 377, R1 S1--S2, mean-speed detector), ICC(A,1) was .375 unfiltered, .727 centred and .320 for the comparator (mean over repetitions; 2.5th--97.5th percentile range of repetitions .135 to .631, not a CI). The contrast centred $-$ random was .407 [.059, .511], and centred $-$ unfiltered was .352 [.020, .526]. At floor 10 ($n$ = 342) the unfiltered ICC was already .727; centred was .759 and the comparator .698. There, centred $-$ random was .061 [.007, .122] and centring (mean\_u $\to$ mean\_c) +.032 [$-$.009, .082] (Supplementary Table S1). The same dependence on the floor held for admission in the joint cohort at each floor (R1 S1--S2): mean\_u $\to$ adm\_u was .308 [.005, .487] at floor 1 ($n$ = 396) and .034 [$-$.008, .082] at floor 10 ($n$ = 331), where mean\_u was already .714 (Supplementary Table S2). The floor dependence also reflected single participants. At floor 10, the participant who carried most of the floor-5 admission contrast (see below) remained eligible for unfiltered mean-speed events but not for any centred construction, and so was outside both the unfiltered/centred intersection and the joint cohort. In the own cohort of unfiltered mean-speed events at floor 10 ($n$ = 351, R1 S1--S2), ICC(A,1) was .371; this participant accounted for 71\% of the summed squared session differences, and omitting them gave .651. The floor dependence of these increments thus coincided with whether a floor admitted a few recordings with extreme $A$.

\subsection{One participant carried most of the admission contrast}

In a sensitivity analysis of unfiltered mean-speed events in their own cohort ($n$ = 379, floor 5, R1 S1--S2), the summed squared difference between the two sessions' $A$ was concentrated in a few participants (\cref{fig:influence}a). The largest contributor accounted for 53\% of the sum and the three largest for 87\%. Each of the two largest contributors had one session whose $A$ was far outside the range of the other participants (1.337 and 1.189; their other sessions, $-$.093 and .068). Removing the largest contributor raised ICC(A,1) from .170 ($n$ = 379) to .381 ($n$ = 378), and removing the two largest gave .646 ($n$ = 377; \cref{fig:influence}b). For centred mean-speed events in their own cohort ($n$ = 377, floor 5, R1 S1--S2), ICC(A,1) was .727; it was between .731 and .737 after removing its one, two, three or five most influential participants, and between .721 and .737 when any single participant was left out. The two constructions in \cref{fig:influence}b use different cohorts, and participants were ranked separately within each construction. Large squared session differences were concentrated in a few participants under the unfiltered construction. The same structure underlay the admission contrast in the joint cohort ($n$ = 376, floor 5, R1 S1--S2; Supplementary Table S13). Under mean\_u, one participant, whose $A$ was 1.189 in session 1 and .068 in session 2 (\cref{fig:influence}d,e), contributed 68\% of the summed squared deviations of the session differences from their mean; under adm\_u the largest contribution was 12\%. The SD of the session differences was .070 under mean\_u and .034 under adm\_u, and the SD of the participants' two-session means was .052 and .040. Omitting this participant from both constructions changed the admission contrast from .319 to .054 [$-$.007, .129] ($n$ = 375; \cref{fig:influence}c): mean\_u rose to .639, whereas adm\_u stayed at .693. Across all 376 single-participant omissions, the contrast ranged from .054 to .338, and no other omission reduced it below .310. The S1--S2 rank correlation was similar under the two constructions (Spearman $\rho$ = .707 for mean\_u and .710 for adm\_u; Supplementary Table S12). Admission therefore changed absolute agreement mainly by removing the events behind one participant's extreme session difference, while the S1--S2 rank correlation changed little. In that participant's first session, the median $l$ of the 22 return sweeps found by mean-speed detection was 2.188, so at least half had a left-eye peak speed about nine or more times the right-eye one. Many of these events combined a left-eye horizontal displacement of more than 40$^\circ$ with little right-eye displacement and left-eye peak speeds above 1,000$^\circ$/s, a level usually treated as physiologically implausible~\citep{nystrom2010}; they were classed as return sweeps because the mean of the two eyes' horizontal amplitudes was large and leftward. Admission kept 11 of the 22 return sweeps and moved their median $l$ to $-$.023 (\cref{fig:influence}f). The worse-eye fixation noise of that session (RMS sample-to-sample deviation .155) was below the median over all recordings (.214), so the data-quality exclusion described below would not have removed it.

% === Fig 3: participant influence (gen_fig3_influence.py) ===
\begin{figure}[t]
\centering
\includegraphics[width=\textwidth]{figures/fig3_influence.pdf}
\caption{\textbf{The low ICC of unfiltered mean-speed events and the admission contrast were carried by a few participants.}
\textbf{a}, Share of the summed squared session differences of $A$, $\Sigma(S1 - S2)^2$ (R1, floor 5), for the ten largest contributors under
unfiltered mean-speed events in their own cohort ($n$ = 379); the top three (black) account for 87\%.
\textbf{b}, ICC(A,1) after removing the $k$ most influential participants (largest absolute leave-one-out shift in ICC), ranked separately within
each mean-speed construction's own cohort ($n$ before removal: unfiltered 379, centred 377); the unfiltered ICC changed from
.170 to .646 after two removals, and the centred ICC stayed between .727 (no removal) and
.737. Because cohorts and rankings differ, a and b do not isolate an effect of centring.
\textbf{c}, Admission contrast (ICC of admitted minus unfiltered mean-speed events) in the joint cohort ($n$ = 376) after each single-participant omission,
against the omitted participant's share of the summed squared deviations of the session differences from their mean under unfiltered mean-speed events.
Dashed line, full-cohort contrast .319 [.003, .501]; star, the omission that gave .054 [$-$.007, .129]
($n$ = 375); brackets, 95\% participant-cluster bootstrap CIs. \textbf{d}, \textbf{e}, Session 2 against session 1 $A$ in the joint cohort for
unfiltered mean-speed (d, ICC = .374) and admitted unfiltered events (e, ICC = .693) on common axes; star, the participant omitted in c;
dotted line, identity. \textbf{f}, For that participant, the median $l = \ln(\mathrm{PV}_L/\mathrm{PV}_R)$ of return sweeps (RS) and of progressive saccades
in each R1 session under the two constructions; numbers, RS counts. Intervals are descriptive. Participant deletion is a sensitivity analysis here, not a cleaning rule.}
\label{fig:influence}
\end{figure}


\subsection{A peak-velocity ceiling, typical session differences and detector settings}

Because the events behind that session were faster than is usually treated as physiologically plausible, we added three post hoc analyses that were not part of the frozen plan (Methods). First, a conventional plausibility comparator removed classified mean-speed events whose larger per-eye peak velocity exceeded 1,000$^\circ$/s, leaving the windows, peaks and classes of the remaining events unchanged. It removed 735 of 138,642 classified mean-speed events (0.53\%; 2.5\% of return sweeps and 0.29\% of progressive saccades). Of the removed events, 40.5\% lacked a matched left/right per-eye pair, whereas only 2.5\% of the 11,952 events without such a pair exceeded the ceiling, so the ceiling and admission removed largely different events. In session 1 of the influential participant, 13 of the 22 return sweeps exceeded the ceiling, and the median $l$ of the remaining ones was $-$.039 (2.188 before); in session 2, 2 of 16 did. All 376 participants of the joint cohort remained eligible. With the ceiling, ICC(A,1) was .711 [.642, .769] (\cref{fig:icc}a). The contrast mean\_u $\to$ ceiling was .337 [.008, .513], and admitted minus ceiling was $-$.018 [$-$.090, .039]; applying admission after the ceiling changed the ICC by .005 [$-$.024, .034] (\cref{fig:icc}b; Supplementary Table S14). With the influential participant omitted ($n$ = 375), mean\_u $\to$ ceiling was .072 [$-$.005, .162] and admitted minus ceiling $-$.018 [$-$.091, .038]; across all 376 single-participant omissions, admitted minus ceiling ranged from $-$.032 to .011. With a 750$^\circ$/s ceiling, admitted minus ceiling was $-$.037 [$-$.108, .014] ($n$ = 374), and with 1,500$^\circ$/s it was .014 [$-$.078, .105] ($n$ = 376). At floor 10 with the 1,000$^\circ$/s ceiling ($n$ = 329), all ten constructions had ICCs between .730 and .762, and none of the contrasts in Supplementary Table S14b exceeded .026 in absolute value. In these data, admission showed no consistent advantage in short-term reliability over a conventional peak-velocity ceiling; the two gave similar point estimates, but equivalence was not tested.

Second, ICC differences describe agreement relative to between-participant spread. To express the session differences in the units of $A$, we computed the absolute R1 S1--S2 difference $|A_{S2} - A_{S1}|$ in the joint cohort (Supplementary Table S15). Its median was .019 [.017, .022] for mean\_u, .017 [.015, .019] for adm\_u and .019 [.016, .021] with the ceiling; the 90th percentiles were .054, .053 and .052. These summary statistics, compared with paired participant-bootstrap draws, were lower with admission: the median by .0022 [.0003, .0041] and the mean by .0047 [.0004, .0118], whereas the difference in the 90th percentile, $-$.0019 [$-$.0071, .0040], included zero. With the influential participant omitted, admission lowered the median by .0020 and the mean by .0017 [.0001, .0036]. Admitted and ceiling constructions differed in median by $-$.0016 [$-$.0030, .0007]. Because $A$ is half a difference of natural-log ratios, a median session difference of .019 corresponds to a multiplicative change of about 4\% ($e^{2 \times .019} \approx 1.04$) in $e^{2A}$, the ratio of the exponentiated class medians of $l$. The median absolute session difference was thus about 11.5\% lower with admission (.0170 against .0192); this compares summary statistics of the cohort and does not describe the change for individual participants.

Third, we re-extracted all R1 recordings with one detector setting changed at a time: velocity threshold 30 or 50$^\circ$/s (frozen, 40), Savitzky--Golay window 5 or 9 samples (frozen, 7) and maximum left/right onset gap for pairing per-eye events 1 or 5 samples (frozen, 3). Re-extraction with the frozen settings reproduced the analysed event sets exactly (Supplementary Table S16). Each variant was evaluated in its own joint cohort ($n$ = 365--377, floor 5), so differences between variants combine the changed setting with changed eligibility; the pairing-gap variants leave the mean-speed events unchanged, and their mean\_u ICCs (.359 and .375, against .374) differ only through the cohort. In these cohorts, the unfiltered mean-speed ICC ranged from .226 (threshold 50$^\circ$/s) to .703 (window 5), whereas admitted ICCs ranged from .661 to .706, ceiling ICCs from .697 to .736 and the four other constructions from .660 to .734 (\cref{fig:icc}d). The admission contrast followed the unfiltered baseline: it was between .282 and .464 in every variant except window 5, and in each of these variants the same participant made the largest contribution to the session differences under mean\_u; omitting them left contrasts of .004 to .054. With window 5, that participant's session-1 $A$ under mean\_u was .160 rather than 1.189, no participant contributed more than 5.4\%, and admission changed the ICC by $-$.029 [$-$.068, .011]. Admitted minus ceiling ranged from $-$.062 to .010 across variants; its interval excluded zero only with window 5 ($-$.062 [$-$.123, $-$.011]), in favour of the ceiling. Whether that session remained extreme under the unfiltered construction, and hence the size of the admission contrast, thus depended on detector settings, whereas the point estimates after any of the plausibility rules stayed between .66 and .74 in the tested variants.

\subsection{Agreement between constructions}

Recording-level $A$ agreed closely between the centred mean-speed and centred per-eye constructions (\cref{fig:agreement}a). Over the 955 recordings from 401 participants eligible under both constructions (floor 5, all rounds), Pearson $r$ was .989 [.984, .992]. The mean difference (per-eye minus mean-speed) was .000, the SD of the differences .007, and the nominal 95\% limits of agreement $-$.013 [$-$.016, $-$.011] to .013 [.011, .015]; 907 of the 955 recordings (95.0\%, 95\% CI 93.5--96.4\%) lay inside the limits, and the empirical 2.5th and 97.5th percentiles of the differences were $-$.013 and .013 (Supplementary Table S13). These statistics describe pooled, recording-level agreement between two constructions, so participants with more eligible recordings carry more weight; they do not describe repeatability between sessions. Comparisons involving unfiltered mean-speed events agreed less well on the absolute scale. Unfiltered per-eye against unfiltered mean-speed (957 recordings, 401 participants) had limits of agreement $-$.077 to .078 (\cref{fig:agreement}b). Because a few large differences inflated the SD, 947 of the 957 recordings (99.0\%) lay inside these limits, and the empirical 2.5th and 97.5th percentiles, $-$.028 and .033, were narrower than the limits. Across the comparisons with unfiltered mean-speed events at floor 5, Pearson $r$ ranged from .746 to .780, whereas Spearman $\rho$ ranged from .956 to .963. Recording ranks thus remained strongly associated across constructions, despite poorer absolute agreement involving unfiltered mean-speed events in a few recordings. Admitted unfiltered and per-eye unfiltered events agreed with $r$ = .985 and limits of agreement of about $\pm$.016.

% === Fig 4: between-construction agreement (gen_fig4_bland_altman.py) ===
\begin{figure}[t]
\centering
\includegraphics[width=\textwidth]{figures/fig4_bland_altman.pdf}
\caption{\textbf{Agreement of recording-level $A$ between constructions.} Bland--Altman plots of all recordings eligible under both
constructions (floor 5, all rounds): vertical axis, per-eye minus mean-speed $A$; horizontal axis, mean of the two. \textbf{a}, Centred per-eye against centred mean-speed construction (955 recordings,
401 participants): bias .0001 [$-$.0004, .0005], nominal 95\% limits of agreement $-$.013 [$-$.016, $-$.011] and .013 [.011, .015];
907 of 955 recordings (95.0\%, 95\% CI 93.5--96.4\%) inside the limits; 2.5th and 97.5th percentiles of the differences $-$.013 [$-$.018, $-$.009] and .013 [.009, .015].
\textbf{b}, Unfiltered per-eye against unfiltered mean-speed construction (957 recordings, 401 participants):
bias .001 [$-$.002, .003], limits $-$.077 [$-$.130, $-$.027] and .078 [.030, .126]; 947 of 957 recordings (99.0\%, 95\% CI 94.5--99.7\%) inside the limits; 2.5th and 97.5th percentiles of the differences $-$.028 [$-$.040, $-$.021] and .033 [.027, .044].
The few large differences widen the limits in b beyond the central spread; the inset shows the recordings with $|$difference$|$ $\leq$ .1. This panel compares detectors at the same filtering.
Solid line, bias; dashed lines, limits of agreement (bias $\pm$ 1.96 SD of the differences); dotted lines, empirical 2.5th and 97.5th percentiles.
All summaries are pooled over recordings, so participants with more eligible recordings carry more weight; brackets are 95\% participant-cluster bootstrap CIs (5,000 draws; Supplementary Table S13).
The recording sets differ between panels, so the panel widths do not isolate the effect of filtering. Pearson $r$ (in-panel text) carries a 95\% participant-cluster bootstrap CI. These are agreement between constructions, not
between-session repeatability.}
\label{fig:agreement}
\end{figure}


\subsection{Repeatability across rounds}

Repeatability over the longer interval depended on aggregation and cohort (\cref{tab:longitudinal}). R3 was recorded about 19--26 months after R1~\citep{lohr2023}. When each round was summarized by the mean of its two sessions, R1--R3 ICC(A,1) in the joint cohort ($n$ = 53, floor 5) ranged from .649 to .709 across the six constructions. The highest value was mean\_c, .709 [.505, .828]. In each unfiltered construction's own cohort ($n$ = 54, floor 5), which added one participant who was ineligible after centring, the round-average ICC(A,1) was lower: .461 [.118, .768] for mean\_u, .468 [.078, .787] for adm\_u and .466 [.084, .761] for pereye\_u. That participant's R3 session-2 recording was ineligible under every centred construction, and all nine of its return sweeps, whose median $l$ was extreme (1.424), had per-eye support, so admission did not remove them. Omitting the participant raised the round-average ICC(A,1) to .649 for mean\_u and .650 for pereye\_u; under mean\_u, they accounted for 46\% of the summed squared differences between round averages. Between single sessions (R1 S1 against R3 S1, joint cohort $n$ = 56, floor 5), ICC(A,1) ranged from .529 [.251, .704] for pereye\_c to .599 [.350, .746] for mean\_u; all six estimates are in \cref{tab:longitudinal}. Round-average and single-session ICCs concern different measurement units, so the columns of \cref{tab:longitudinal} are not directly comparable, and their differences do not identify an effect of elapsed time. Longitudinal contrasts varied across constructions, aggregation schemes and cohorts. In the joint longitudinal cohort, centred per-eye $A$ had a lower round-average ICC than centred mean-speed $A$ (mean\_c $\to$ pereye\_c $\Delta$ICC = $-$.059 [$-$.137, $-$.008], $n$ = 53). For R1--R2 round averages, centred mean-speed $A$ had ICC(A,1) .667 [.338, .838] ($n$ = 29, floor 5).

% === Table 1: longitudinal repeatability (gen_table1_longitudinal.py) ===
\input{figures/TABLE_1.tex}


\subsection{What recording-centred abstention detects}

Because centring rejects events whose $l$ is unusual for the recording, we injected synthetic one-eye faults into R1 recordings and asked how often centring, applied to mean-speed events, abstained from the injected events (\cref{fig:injection}; Supplementary Table S9). Admission and per-eye detection were not evaluated by these injections. Recall was computed among injections reaching centring, that is, injections overlapped by a mean-speed event that passed the detection, class and exclusion rules. In this subsection, ``reaching centring'' is distinct from per-eye admission. The frozen criterion for judged conditions was recall of at least .80 with a participant-cluster lower bound above .70. Tested one-eye steps (pseudo-saccades in one eye during fixation) met it. Conditional on events reaching centring that overlapped the injected onset, centring rejected 4,771 of 4,780 4$^\circ$ steps (recall .998 [.997, .999]) and all 4,780 8$^\circ$ steps; because every observed injection was rejected, every participant-bootstrap draw gave a recall of 1, and this degenerate interval does not quantify uncertainty about unobserved failures. The 2$^\circ$ step (4,198 of 4,308; .974 [.969, .980]) was not judged. The rule did not meet the frozen criterion for one-eye slowing by a factor of 3 (1,881 of 2,437; recall .772 [.754, .789]; .504 of injections reached centring) or for a 6$^\circ$ peak glitch in one eye (1,807 of 4,803; .376 [.363, .389]). One-eye dropouts were mostly removed before centring: 98.5\% of dropout injections were excluded at the missingness layer, and only 61 of 4,835 reached centring, so their recall was not interpreted. As a negative control, mild one-eye displacement scaling increased rejection of the same target events by at most .009 (maximum of the observed point estimates; upper bounds in Supplementary Table S9). The keep/abstain status of untouched events changed in at most 0.29\% of events (maximum observed point estimate; upper bound 0.35\%). In the contamination conditions (5--30\% of eligible events injected), at most 2.8\% of injections overlapped another injection, and recall on non-overlapping injections was within .0005 of the overall recall. Overlap was more frequent in the mixed-contamination condition described below (47 of 686 injected slowings, 6.9\%, and 31 of 45 injected 4$^\circ$ steps); recall on its non-overlapping events was .803, against .799 overall (Supplementary Table S7).

% === Fig 5: synthetic one-eye faults (gen_fig5_injection.py) ===
\begin{figure}[t]
\centering
\includegraphics[width=\textwidth]{figures/fig5_injection.pdf}
\caption{\textbf{Abstention by the centring rule under synthetic one-eye faults.}
\textbf{a}, Recall: fraction of injections reaching centring (overlapped by a detected event that passed the class and exclusion rules) that the centring rule abstained from, with
95\% participant-cluster bootstrap CI; denominator, injections reaching centring. Dashes show the fraction reaching centring; denominator, all injections. The frozen criterion for judged conditions (filled) was
a point estimate of at least .80 (dashed line) with lower bound $>$ .70; open blue markers are unjudged magnitudes, grey markers are descriptive contamination
conditions (one-eye slowing by a factor of 3 applied to 5--30\% of eligible events). Of the judged one-eye conditions, only the tested 4$^\circ$ and 8$^\circ$ steps met the criterion; one-eye slowing (factor 3) and a
6$^\circ$ peak glitch did not. For one-eye dropout, only 61 of 4,835 injections reached centring (most were removed at the missingness layer),
so its recall was not interpreted. \textbf{b}, Specificity control: the within-saccade displacement of one randomly chosen eye was scaled by
$\kappa$, a negative-control perturbation ($|\Delta l| \approx .15 < \log 2$). Excess rejection is the rejection rate of the same
target events reaching centring after scaling minus that in the uninjected baseline pass ($\kappa$ = 0.85: .009 [.006, .012];
$\kappa$ = 1.15: .007 [.005, .010]; participant-cluster bootstrap).
Injection counts, the fraction reaching centring, recall or excess rejection, untouched-event changes and outcome layers per condition are in Supplementary Table S9.}
\label{fig:injection}
\end{figure}


On uncontaminated data, events rejected by centring mostly lacked per-eye support: 73.0\% of rejected mean-speed events (8,945 of 12,256) overlapped no matched per-eye pair, against 2.4\% of kept events (3,007 of 126,386; odds ratio 110.8 [97.9, 125.2]; Cohen's $\kappa$ = .714 [.702, .726]). The fraction of a recording's events rejected by centring correlated with the worse eye's fixation noise (RMS sample-to-sample deviation; Spearman $\rho$ = .545 [.487, .600]; 1,004 recordings, 407 participants; Supplementary Table S10). Excluding participants with any recording above the 90th percentile of worse-eye fixation noise left the centred mean-speed ICC(A,1) at .739 [.682, .788] ($n$ = 325, floor 5, R1 S1--S2).

In a mixed synthetic contamination, evaluated on the common cohort of its four runs ($n$ = 376, whose membership differs from the six-construction joint cohort), one R1 session of 20 participants received one-eye slowing on 30\% of eligible events and 4$^\circ$ one-eye steps in 10\% of fixations. The centred mean-speed ICC(A,1) over the full sample changed from .726 to .722 ($\Delta$ = $-$.004 [$-$.009, $-$.0004]; $n$ = 376, floor 5, R1 S1--S2). On the contaminated data, centred minus unfiltered was .341 [.007, .515]. In the 17 injected participants with both sessions eligible in all runs, centred ICC(A,1) went from .478 to .338 ($\Delta$ = $-$.140 [$-$.412, .031]) and unfiltered from .465 to .283 ($\Delta$ = $-$.182 [$-$.273, $-$.014]). Two of the injected participants lost centred eligibility at floor 5 and one lost unfiltered eligibility (Supplementary Table S6).

\subsection{Amplitude-matched left/right controls}

$A$ contrasts saccades that differ in direction, amplitude and role in reading. We therefore also paired leftward and rightward saccades matched on amplitude and position within each recording and computed $D = \tfrac12(l_{\mathrm{leftward}} - l_{\mathrm{rightward}})$ (\cref{fig:matched}; Supplementary Table S8). Participant $D$ is the mean of the two R1 session medians. It was positive in the reading task (TEX; mean .020 [.011, .029], $n$ = 351, 5,617 pairs, 68.7\% of participants above zero) and in the random-saccade task (RAN; .039 [.028, .050], $n$ = 114). Participant TEX $D$ correlated with the two-session mean of TEX $A$ (R1, centred mean-speed construction; Spearman $\rho$ = .446 [.348, .536], $n$ = 331), and so did RAN $D$ ($\rho$ = .483 [.307, .635], $n$ = 112). The TEX association did not require shared events. With $A$ from session 1 and $D$ from session 2, $\rho$ was .398 [.276, .505] ($n$ = 239); with $A$ from session 2 and $D$ from session 1, $\rho$ was .391 [.270, .501] ($n$ = 246). Same-session minus cross-session $\rho$ was $-$.024 [$-$.100, .049] ($n$ = 185). TEX $D$ itself had a low repeatability between R1 sessions (ICC(A,1) = .319 [.159, .472], $n$ = 193, participants with at least 5 pairs per session). RAN $D$ was thus positively associated with reading $A$, but its test--retest reliability could not be estimated, because no participant had at least 5 matched pairs in each R1 session. These controls agree in sign with an abducting eye that is faster~\citep{collewijn1988}. Because direction, amplitude range and event class co-vary within $A$, they do not identify the physiological source of $A$.

% === Fig 6: amplitude-matched peak-velocity difference (gen_fig6_matched.py) ===
\begin{figure}[t]
\centering
\includegraphics[width=\textwidth]{figures/fig6_matched.pdf}
\caption{\textbf{Amplitude-matched left/right peak-velocity difference $D$ and its association with $A$.}
$D = \tfrac12(l_{\mathrm{leftward}} - l_{\mathrm{rightward}})$ for a leftward and a rightward saccade matched on amplitude and position within a
recording (dimensionless, log units); positive $D$ means a relatively faster left eye in leftward than in rightward saccades. Participant $D$ requires
both R1 sessions and at least 5 matched pairs in total; its repeatability requires at least 5 pairs in each session, hence the smaller $n$ below.
\textbf{a}, Participant $D$ (mean of the two R1 session medians) in the reading (TEX; $n$ = 351, mean .020 [.011, .029])
and random-saccade tasks (RAN; $n$ = 114, mean .039 [.028, .050]); histograms are normalized to density; values beyond the axis range are placed in
the edge bins and counted in the legend. Different TEX and RAN sample sizes do not by themselves establish a task difference. \textbf{b}, Participant $D$ against TEX $A$ (R1 two-session mean, centred mean-speed construction):
TEX $D$, Spearman $\rho$ = .446 [.348, .536] ($n$ = 331); RAN $D$, $\rho$ = .483 [.307, .635]
($n$ = 112); $D$ values below $-$.3 are drawn as triangles at the lower edge, and their number is given in the panel. \textbf{c}, Spearman $\rho$ with 95\% participant-bootstrap CI between single-session TEX $A$ and single-session $D$ (same- and cross-session pairings,
$n$ = 239--255) and, below the line, for the pooled associations of panel b; cohorts and averaging differ between the two groups, so their magnitudes are
compared descriptively. TEX $D$ repeatability (R1 S1--S2):
ICC(A,1) = .319 [.159, .472], $n$ = 193. RAN repeatability was not estimable, because no participant had at least 5 pairs in each session, and the RAN session-specific associations
rest on small samples (Supplementary Table S8).}
\label{fig:matched}
\end{figure}


\subsection{Robustness to constant per-eye gain}

For a fixed set of events and classes, a constant gain applied to one eye's signal shifts every $l$ in a recording by the same amount, which cancels in $A$ (Methods). Gain applied to the raw gaze signal can, however, change which events are detected. We therefore re-extracted all events after applying constant affine per-eye gains to the raw gaze (centred mean-speed construction, floor 5). All 11 judged conditions passed the pre-specified criteria (Supplementary Table S11). The largest median $|\Delta A|$ was .0017, the largest 95th percentile of $|\Delta A|$ was .0125, and the largest $|\Delta\mathrm{ICC}|$ was .0068. Eligible coverage changed by $-$.30 to +.20 percentage points. The six exploratory conditions with an absolute log interocular gain ratio $|\delta|$ = .30 also passed.
